\documentclass[11pt]{article}
\usepackage{amsmath,amssymb,geometry}
\geometry{margin=27mm}
\title{EPPA--III: Claim 3.4 block-system audit}
\date{10 October 2026}
\begin{document}
\maketitle
\paragraph{Source anchor.}
The first unnumbered claim of Lemma~\texttt{lemma:doublecover},
source lines 947--964 in \texttt{main-new(7).tex},
asserts that the nontrivial invariant partition of the
\(2n\)-vertex host has \(n\) blocks of size two,
and \(V(G)\) meets each block in precisely one vertex.
\paragraph{Two-point orbit reduction.}
Every isomorphism of ordered edges (and separately of
ordered nonedges) of \(G\) extends to an automorphism of
\(H\) which preserves the block relation. Thus each
adjacency type of \(G\) has a uniform within-block
or cross-block status. Lean verifies this in
\texttt{EPPAIII/BlockSystems/Basic.lean}.
\paragraph{Excluding the three unwanted cases.}
If both types are internal, then \(G\) lies in one block.
All blocks have equal size \(m\le n\), and that block
contains \(n\) vertices of \(G\), so it equals \(V(G)\).
An EPPA extension of any nonempty partial automorphism
of \(G\) stabilizes the block, hence restricts to
an automorphism of \(G\); the empty map extends trivially.
Thus \(G\) would be homogeneous.
If edges are internal and nonedges cross, \(G\) is
a disjoint union of cliques. If nonedges are internal
and edges cross, \(G\) is a complement of such a union.
Both are excluded by the lemma hypotheses.
\paragraph{Transversality via induced triples.}
The exclusion of disjoint unions of cliques is
equivalent to an induced \(P_3\), and its complement
to an induced complement of \(P_3\).
These two triples force distinct selected vertices
to occupy distinct blocks, unless all \(G\) vertices
are in one block. This intermediate implication is
formalized in \texttt{BlockSystems/Transversal.lean}.
\paragraph{Counting.}
Let \(b\) denote the number of blocks and
\(m\ge2\) their common size. Transversality gives
\(b\ge n\). Therefore \(2n=bm\ge2b\ge2n\),
so \(b=n\) and \(m=2\). The arithmetic implication
is formalized in
\texttt{BlockSystems.transversal\_double\_count}.
\paragraph{C16: suggested one-sentence clarification.}
The source proof compresses the case with all of
\(G\) in one block into the assertion that
this block is an EPPA witness. To justify it
without changing the proof structure, insert:
\begin{quote}
Indeed, \(\Pi\) contains \(V(G)\) and
\(|\Pi|\le n\), so \(\Pi=V(G)\);
an extension of a nonempty partial automorphism
of \(G\) stabilizes \(\Pi\) and thus restricts
to an automorphism of \(G\).
\end{quote}
\paragraph{Updated formalization checkpoint.}
The chain of Lean modules
\texttt{Basic.lean},
\texttt{Transversal.lean},
\texttt{ClusterObstructions.lean},
\texttt{FilledBlock.lean},
\texttt{Size.lean} and
\texttt{TransitivePartition.lean}
proves Claim~3.4's invariant two-point partition
directly from vertex-transitivity, EPPA,
the nontrivial imprimitivity system, and
the graph-class exclusions.
\texttt{CanonicalPairs.lean}
constructs the equivalence
\(W\simeq\{0,1\}\times V(G)\) sending
the bottom copy to the original embedding.
\texttt{GraphTransport.lean} proves
that the EPPA and invariant-block conditions
transport through this equivalence.
Finally,
\texttt{FullDoubleCover.lean}
establishes the structural Taylor-double
conclusion with optional mate matching.
All these theorems compile and pass the
permitted-axiom check.

\paragraph{Remaining scope.}
The manuscript additionally asserts transitivity
of \(\operatorname{Aut}(H)\) on ordered inter-block
edges and nonedges, which is \emph{not}
contained in the present structural theorem.
Using the manuscript's exact exclusion of
\(C_5\), \(R_3\), and clique unions to obtain
the required non-homogeneity relies on
the Gardiner classification in its introduction,
not on a Lean proof of that classification.
No original manuscript TeX has been modified.
\end{document}
